\cleardoublepage
\section{辐射}
\subsection{偶极辐射}
\subsubsection{电偶极子辐射}
假设两个小金属球相距为$d$并用一个细导线相连。在时刻$t$，上面的小球带电$q(t)$，下面的为$-q(t)$。假设以角频率$\omega$驱动电荷通过导线在两端的小球上来回振荡，$q(t)=q_0\cos(\omega t)$这样产生一个振荡的电偶极子：
\begin{equation}
    p(t)=p_0\cos(\omega t)\boldsymbol{e}_z
\end{equation}
推迟势为
\begin{equation}
    U=\frac{1}{4\pi\varepsilon_0}\left\{\frac{q_0\cos\left[\omega\left(t-n_+/c\right)\right]}{n_+}
    -\frac{q_0\cos\left[\omega\left(t-n_-/c\right)\right]}{n_-}\right\}
\end{equation}
式中:
$$n_\pm =\sqrt{r^2\mp rd\cos\theta+\left(d/2\right)^2}$$
现在把物理上的电偶极子变成理想的电偶极子，使正负电荷相距非常小$d<<r$：
\begin{align*}
    n_\pm &\cong   r\left(1\mp\frac{d}{2r}\cos\theta\right)\\
    \cos\left[\omega\left(t-n_\pm/c\right)\right]&=
    \cos\left[\omega\left(t-\frac{r}{c}\mp\frac{d}{2c}\cos\theta\right)\right]\\
    \cos\left[\omega\left(t-n_\pm/c\right)\right]&=
    \cos\left[\omega\left(t-\frac{r}{c}\right)\right]\cos\left(\frac{\omega d}{2c}\cos\theta\right)\mp
    \sin\left[\omega\left(t-\frac{r}{c}\right)\right]\sin\left(\frac{\omega d}{2c}\cos\theta\right)
\end{align*}
假设$d<<\dfrac{c}{\omega}$
\begin{align*}
    \cos\left[\omega\left(t-n_\pm/c\right)\right]&\cong
    \cos\left[\omega\left(t-\frac{r}{c}\right)\right]\mp
    \sin\left[\omega\left(t-\frac{r}{c}\right)\right]\frac{\omega d}{2c}\cos\theta\\
    U&=\frac{q_0\cos\left[\omega\left(t-\frac{r}{c}\right)\right]}{4\pi\varepsilon_0}\left\{\frac{1-
    \frac{\omega d}{2c}\cos\theta\tan\left[\omega\left(t-\frac{r}{c}\right)\right]}{r-\frac{d}{2}\cos\theta}
    -\frac{1+
    \frac{\omega d}{2c}\cos\theta\tan\left[\omega\left(t-\frac{r}{c}\right)\right]}{r+\frac{d}{2}\cos\theta}\right\}\\
    U&=\frac{q_0\cos\left[\omega\left(t-\frac{r}{c}\right)\right]\cos\theta}{4\pi\varepsilon_0}\left\{\frac{
    4d-4\frac{r\omega d}{c}\tan\left[\omega\left(t-\frac{r}{c}\right)\right]}{4r^2-d^2\cos^2\theta}
    \right\}\\
    U&=\frac{q_0d\cos\left[\omega\left(t-\frac{r}{c}\right)\right]\cos\theta}{4\pi\varepsilon_0}\left\{\frac{
    4c-4r\omega\tan\left[\omega\left(t-\frac{r}{c}\right)\right] }{4r^2c}
    \right\}\\
\end{align*}
假设$r>>\dfrac{c}{\omega}$
\begin{equation}
    U=-\frac{q_0d\sin\left[\omega\left(t-\frac{r}{c}\right)\right]}{4\pi\varepsilon_0}\frac{
    \omega \cos\theta}{rc}
\end{equation}
\begin{align}
    I&=\frac{\dif q}{\dif t}\nonumber\\
    I&=-q_0\omega\sin(\omega t)\nonumber\\
    \boldsymbol{A}&=\frac{\mu_0}{4\pi}\int_{\frac{d}{2}}^{\frac{d}{2}}\frac{-q_0\omega\sin\left(t\omega-\frac{r\omega}{c}\right)}{r}\boldsymbol{e}_z\dif z\nonumber\\
    \boldsymbol{A}&=\frac{\mu_0}{4\pi}\frac{-d_0\omega\sin\left(t\omega-\frac{r\omega}{c}\right)}{r}\boldsymbol{e}_z
\end{align}
\begin{align}
    \boldsymbol{\nabla}U&=\frac{\partial U}{\partial r}\boldsymbol{e}_r+\frac{\partial U}{r\partial \theta}\boldsymbol{e}_\theta\nonumber\\
    \frac{\partial U}{\partial r}&=-\frac{\partial \frac{p_0\sin\left[\omega\left(t-\frac{r}{c}\right)\right]}{4\pi\varepsilon_0}\frac{
    \omega \cos\theta}{rc}}{\partial r}\nonumber\\
    \frac{\partial U}{\partial r}&=\frac{p_0\sin\left[\omega\left(t-\frac{r}{c}\right)\right]}{4\pi\varepsilon_0}\frac{
    \omega \cos\theta}{r^2c}
    +\frac{\omega^2 \cos\theta}{rc^2}\frac{p_0\cos\left[\omega\left(t-\frac{r}{c}\right)\right]}{4\pi\varepsilon_0}\nonumber\\
    \frac{\partial U}{\partial r}&\cong\frac{\omega^2 \cos\theta}{rc^2}\frac{p_0\cos\left[\omega\left(t-\frac{r}{c}\right)\right]}{4\pi\varepsilon_0}\nonumber\\
    \frac{\partial U}{r\partial \theta}&=-\frac{\partial \frac{q_0d\sin\left[\omega\left(t-\frac{r}{c}\right)\right]}{4\pi\varepsilon_0}\frac{
    \omega \cos\theta}{rc}}{r\partial \theta}\nonumber\\
    \frac{\partial U}{r\partial \theta}&=\frac{q_0d\sin\left[\omega\left(t-\frac{r}{c}\right)\right]}{4\pi\varepsilon_0}\frac{
    \omega \sin\theta}{r^2c}\nonumber\\
    \boldsymbol{\nabla}U&\cong\frac{\omega^2 \cos\theta}{rc^2}\frac{p_0\cos\left[\omega\left(t-\frac{r}{c}\right)\right]}{4\pi\varepsilon_0}\boldsymbol{e}_r\\
    \frac{\partial \boldsymbol{A}}{\partial t}&=\frac{\partial \frac{\mu_0}{4\pi}\frac{-p_0\omega\sin\left(t\omega-\frac{r\omega}{c}\right)}{r}\boldsymbol{e}_z}{\partial t}\nonumber\\
    \frac{\partial \boldsymbol{A}}{\partial t}&=-\frac{\mu_0}{4\pi}\frac{p_0\omega^2\cos\left(t\omega-\frac{r\omega}{c}\right)}{r}\boldsymbol{e}_z\nonumber\\
    \frac{\partial \boldsymbol{A}}{\partial t}&=-\frac{p_0\omega^2\cos\left(t\omega-\frac{r\omega}{c}\right)}{4\pi\varepsilon_0 rc^2}\left(\cos\theta\boldsymbol{e}_r-\sin\theta\boldsymbol{e}_\theta\right)\nonumber\\
    \boldsymbol{E}&=-\boldsymbol{\nabla}U-\frac{\partial \boldsymbol{A}}{\partial t}\nonumber\\
    \boldsymbol{E}&=-\frac{p_0\omega^2\cos\left(t\omega-\frac{r\omega}{c}\right)}{4\pi\varepsilon_0 rc^2}\sin\theta\boldsymbol{e}_\theta\\
    \boldsymbol{B}&=\boldsymbol{\nabla}\times\boldsymbol{A}\nonumber\\
    \boldsymbol{B}&=\boldsymbol{\nabla}\times\frac{\mu_0}{4\pi}\frac{-p_0\omega\sin\left(t\omega-\frac{r\omega}{c}\right)}{r}\left(\cos\theta\boldsymbol{e}_r-\sin\theta\boldsymbol{e}_\theta\right)\nonumber\\
    \boldsymbol{B}&=\frac{1}{r}\left(\frac{\partial ra_\theta}{\partial r}
    -\frac{\partial a_r}{\partial \theta}\right)\boldsymbol{e}_\phi\nonumber\\
    \boldsymbol{B}&=\frac{\mu_0p_0\omega}{4\pi r}\left[\frac{\partial \sin\left(t\omega-\frac{r\omega}{c}\right)\sin\theta}{\partial r}
    +\frac{\partial \sin\left(t\omega-\frac{r\omega}{c}\right)\cos\theta}{r\partial \theta}\right]\boldsymbol{e}_\phi\nonumber\\
    \boldsymbol{B}&=-\frac{\mu_0p_0\omega}{4\pi r}\left[\frac{\omega \cos\left(t\omega-\frac{r\omega}{c}\right)\sin\theta}{c}
    +\frac{\sin\left(t\omega-\frac{r\omega}{c}\right)\sin\theta}{r}\right]\boldsymbol{e}_\phi\nonumber\\
    \boldsymbol{B}&\cong-\frac{\mu_0p_0\omega^2}{4\pi rc}\cos\left(t\omega-\frac{r\omega}{c}\right)\sin\theta\boldsymbol{e}_\phi\\
    \boldsymbol{S}&=\frac{1}{\mu_0}\boldsymbol{E}\times\boldsymbol{B}\nonumber\\
    \boldsymbol{S}&=\frac{1}{\mu_0}\frac{d_0\omega^2\cos\left(t\omega-\frac{r\omega}{c}\right)}{4\pi\varepsilon_0 rc^2}\sin\theta\boldsymbol{e}_\theta\times\frac{\mu_0d_0\omega^2}{4\pi rc}\cos\left(t\omega-\frac{r\omega}{c}\right)\sin\theta\boldsymbol{e}_\phi\nonumber\\
    \boldsymbol{S}&=\frac{d_0^2\omega^4}{16\pi^2\varepsilon_0 r^2c^3}\sin^2\theta\cos^2\left(t\omega-\frac{r\omega}{c}\right)\boldsymbol{e}_r\\
    \langle\boldsymbol{S} \rangle&=\frac{d_0^2\omega^4}{32\pi^2\varepsilon_0 r^2c^3}\sin^2\theta\boldsymbol{e}_r\\
    \langle P\rangle&=\int\boldsymbol{S}\cdot\dif\boldsymbol{a}\nonumber\\
    \langle P\rangle&=\frac{d_0^2\omega^4}{12\pi\varepsilon_0c^3}
\end{align}
\subsubsection{磁偶极子辐射}
假设有一个由导线构成的半径为$b$的环，它通有一个交变电流$I=I_0\cos\omega t$，这就是一个振动的磁偶极子模型$\boldsymbol{m}=\pi b^2I\boldsymbol{e}_z$
\begin{align}
    \boldsymbol{A}&=\frac{\mu_0}{4\pi}\int\frac{I_0\cos\left(t\omega-\frac{n\omega}{c}\right)}{n}\dif \boldsymbol{l}\nonumber
\end{align}
对于一个在$x$轴上方的点$r$，$\boldsymbol{A}$必定指向$y$方向，因为对称分布在$x$轴两边的$x$分量彼此抵消，故
\begin{align}
    \boldsymbol{A}&=\frac{\mu_0b}{4\pi}\int_0^{2\pi}\frac{I_0\cos\left(t\omega-\frac{n\omega}{c}\right)}{n}\cos\phi\boldsymbol{e}_y\dif \phi\nonumber
\end{align}
假设1:$b<<r$：
\begin{align*}
    n&\cong r\left(1-\frac{b}{r}\sin\theta\cos\phi\right)\\
    \cos\left(t\omega-\frac{n\omega}{c}\right)&\cong\cos\left[t\omega-\frac{r\left(1-\frac{b}{r}\sin\theta\cos\phi\right)\omega}{c}\right]\\
    \cos\left(t\omega-\frac{n\omega}{c}\right)&=\cos\left(t\omega-\frac{r\omega}{c}\right)\cos\left(\frac{b\sin\theta\cos\phi\omega}{c}\right)
    -\sin\left(t\omega-\frac{r\omega}{c}\right)\sin\left(\frac{b\sin\theta\cos\phi\omega}{c}\right)
\end{align*}
假设2:$b<<\dfrac{c}{\omega}$
\begin{align*}
    \cos\left(t\omega-\frac{n\omega}{c}\right)&=\cos\left(t\omega-\frac{r\omega}{c}\right)
    -\sin\left(t\omega-\frac{r\omega}{c}\right)\frac{b\sin\theta\cos\phi\omega}{c}\\
    \boldsymbol{A}&\cong\frac{\mu_0I_0}{4\pi}\int_0^{2\pi}\frac{\cos\left(t\omega-\frac{r\omega}{c}\right)
    -\sin\left(t\omega-\frac{r\omega}{c}\right)\frac{b\sin\theta\cos\phi\omega}{c}}{r\left(1-\frac{b}{r}\sin\theta\cos\phi\right)}\cos\phi\boldsymbol{e}_y\dif \phi\\
    A&=\frac{\mu_0I_0b}{4\pi }\int_0^{2\pi}\frac{
    c\cos\left(t\omega-\frac{r\omega}{c}\right)-\sin\left(t\omega-\frac{r\omega}{c}\right)b\sin\theta\cos\phi\omega}{cr^2}\left(r+b\sin\theta\cos\phi\right)\cos\phi\dif \phi\\
    A&\cong\frac{\mu_0I_0b}{4\pi }\int_0^{2\pi}\frac{
    c\cos\left(t\omega-\frac{r\omega}{c}\right)b\sin\theta\cos\phi-\sin\left(t\omega-\frac{r\omega}{c}\right)b\sin\theta\cos\phi\omega r}{cr^2}\cos\phi\dif \phi\\
    A&=\frac{\mu_0I_0b^2}{4}\frac{
    c\cos\left(t\omega-\frac{r\omega}{c}\right)-\sin\left(t\omega-\frac{r\omega}{c}\right)\omega r}{cr^2}\sin\theta
\end{align*}
假设3:$r>>\dfrac{c}{\omega}$
\begin{align}
    \boldsymbol{A}&=-\frac{\mu_0m_0}{4\pi}\frac{
    \sin\left(t\omega-\frac{r\omega}{c}\right)\omega }{cr}\sin\theta\boldsymbol{e}_\phi\\
    \boldsymbol{E}&=\frac{\mu_0m_0}{4\pi}\frac{
    \cos\left(t\omega-\frac{r\omega}{c}\right)\omega^2 }{cr}\sin\theta\boldsymbol{e}_\phi\\
    \boldsymbol{B}&=-\frac{\mu_0m_0}{4\pi}\frac{
    \cos\left(t\omega-\frac{r\omega}{c^2}\right)\omega^2 }{cr}\sin\theta\boldsymbol{e}_\theta\\
    \boldsymbol{S}&=\frac{\mu_0}{c}\left[\frac{m_0}{4\pi}\frac{
    \cos\left(t\omega-\frac{r\omega}{c}\right)\omega^2 }{cr}\sin\theta\right]^2\boldsymbol{e}_r\\
    \langle\boldsymbol{S} \rangle&=\frac{\mu_0}{c}\frac{m_0^2}{32\pi^2}\frac{\omega^4 }{c^2r^2}\sin^2\theta\boldsymbol{e}_r\\
    \langle P\rangle&=\frac{\mu_0m_0^2\omega^4}{12\pi c^3}
\end{align}
\subsubsection{任意源的辐射}
假设1:$r'<<r$：
\begin{align*}
    n&\cong r\left(1-\frac{\boldsymbol{r}\cdot\boldsymbol{r}'}{r^2}\right)\\
    \frac{1}{n}&\cong \frac{1}{r}\left(1+\frac{\boldsymbol{r}\cdot\boldsymbol{r}'}{r^2}\right)\\
    \rho\left(\boldsymbol{r}',t-\frac{n}{c}\right)&\cong\rho\left(\boldsymbol{r}',t-\frac{r}{c}+\frac{\boldsymbol{r}\cdot\boldsymbol{r}'}{rc }\right)
\end{align*}
定义$t_0\equiv t-\frac{n}{c}$泰勒展开得
\begin{align*}
    \rho\left(\boldsymbol{r}',t-\frac{n}{c}\right)=\rho\left(\boldsymbol{r}',t_0\right)
    +\dot\rho\left(\boldsymbol{r}',t_0\right)\frac{\boldsymbol{r}\cdot\boldsymbol{r}'}{rc }+...
\end{align*}
假设2:泰勒展开的高阶项可略去：
\begin{align*}
    U&=\frac{1}{4\pi\varepsilon_0}\int\frac{\rho}{n}\dif\tau\\
    U&\cong\frac{1}{4\pi\varepsilon_0}\int\frac{\rho\left(\boldsymbol{r}',t_0\right)
    +\dot\rho\left(\boldsymbol{r}',t_0\right)\frac{\boldsymbol{r}\cdot\boldsymbol{r}'}{rc }}{r}\left(1+\frac{\boldsymbol{r}\cdot\boldsymbol{r}'}{r^2}\right)\dif\tau\\
    U&\cong\frac{1}{4\pi\varepsilon_0r}\left(\int \rho\dif\tau
    +\frac{\boldsymbol{r}}{r^2}\cdot\int\boldsymbol{r'} \rho\dif\tau
    +\frac{\boldsymbol{r}}{r c}\cdot\frac{\dif}{\dif t}\int \boldsymbol{r}'\rho\dif\tau\right)\\
    U&\cong\frac{1}{4\pi\varepsilon_0r}\left[Q
    +\frac{\boldsymbol{r}\cdot\boldsymbol{d}(t_0)}{r^2}\cdot
    +\frac{\boldsymbol{r}\cdot\dot{\boldsymbol{d}}(t_0)}{r c}\right]
\end{align*}
\begin{align*}
    \boldsymbol{A}&=\frac{\mu_0}{4\pi}\int\frac{\boldsymbol{J}(\boldsymbol{r}',t_0)}{n}\dif\tau\\
    \boldsymbol{A}&\cong\frac{\mu_0}{4\pi}\int\frac{\boldsymbol{J}(\boldsymbol{r}',t_0)}{r}\dif\tau\\
    \boldsymbol{A}&\cong\frac{\mu_0}{4\pi}\frac{\dot{\boldsymbol{d}}(t_0)}{r}\\
    \boldsymbol{\nabla}t_0&=-\frac{1}{c}\boldsymbol{\nabla}r\\
    \boldsymbol{\nabla}t_0&=-\frac{\boldsymbol{e}_r}{c}
\end{align*}
因为求的是辐射，所以不对变量$t_0$求导的项对辐射没有贡献
\begin{align*}
    \boldsymbol{\nabla}U&\cong\boldsymbol{\nabla}\left[\frac{1}{4\pi\varepsilon_0r}\frac{\boldsymbol{r}\cdot\dot{\boldsymbol{d}}(t_0)}{r c}\right]\\
    \boldsymbol{\nabla}U&\cong\frac{1}{4\pi\varepsilon_0}\frac{\boldsymbol{r}\cdot\ddot{\boldsymbol{d}}(t_0)}{r^2 c}\boldsymbol{\nabla}t_0\\
    \boldsymbol{\nabla}U&\cong-\frac{1}{4\pi\varepsilon_0}\frac{\boldsymbol{r}\cdot\ddot{\boldsymbol{d}}(t_0)}{r^2 c^2}\boldsymbol{e}_r\\
    \boldsymbol{\nabla}\times\boldsymbol{A}&\cong\frac{\mu_0}{4\pi r}\boldsymbol{\nabla}\times\dot{\boldsymbol{d}}(t_0)\\
    \boldsymbol{\nabla}\times\boldsymbol{A}&\cong\frac{\mu_0}{4\pi r}\left(\boldsymbol{\nabla}t_0\right)\times\ddot{\boldsymbol{d}}(t_0)\\
    \boldsymbol{\nabla}\times\boldsymbol{A}&\cong-\frac{\mu_0}{4\pi rc}\boldsymbol{e}_r\times\ddot{\boldsymbol{d}}(t_0)\\
    \frac{\partial \boldsymbol{A}}{\partial t}&\cong\frac{\mu_0}{4\pi r}\ddot{\boldsymbol{d}}(t_0)
\end{align*}
\begin{align}
    \boldsymbol{E}(\boldsymbol{r},t)&\cong-\boldsymbol{\nabla}U-\frac{\partial \boldsymbol{A}}{\partial t}\nonumber\\
    \boldsymbol{E}(\boldsymbol{r},t)&\cong\frac{\mu_0}{4\pi r}\left\{\left[\boldsymbol{e}_r\cdot\ddot{\boldsymbol{d}(t_0)}\right]\boldsymbol{e}_r-\ddot{\boldsymbol{d}}(t_0)\right\}\nonumber\\
    \boldsymbol{E}(\boldsymbol{r},t)&\cong\frac{\mu_0}{4\pi r}\left\{\boldsymbol{e}_r\times\left[\boldsymbol{e}_r\times\ddot{\boldsymbol{d}(t_0)}\right]\right\}\\
    \boldsymbol{B}&=\boldsymbol{\nabla}\times\boldsymbol{A}\nonumber\\
    \boldsymbol{B}&\cong-\frac{\mu_0}{4\pi rc}\boldsymbol{e}_r\times\ddot{\boldsymbol{d}}(t_0)
\end{align}
如果使用球坐标，令$z$轴指向$\ddot{\boldsymbol{d}}(t_0)$方向
\begin{align}
    \boldsymbol{E}&\cong\frac{\mu_0\sin\theta}{4\pi r}\ddot{\boldsymbol{d}}(t_0)\boldsymbol{e}_\theta\\
    \boldsymbol{B}&\cong\frac{\mu_0\sin\theta}{4\pi rc}\ddot{\boldsymbol{d}}(t_0)\boldsymbol{e}_\phi\\
    \boldsymbol{S}&\cong\frac{\mu_0\sin^2\theta}{16\pi^2 r^2c}\ddot{\boldsymbol{d}}^2(t_0)\boldsymbol{e}_r\\
    P&\cong\frac{\mu_0\ddot{\boldsymbol{d}}^2(t_0)}{6\pi c}
\end{align}
\mysssec{求在偶极子中连接两端导线的“辐射电阻”}
\begin{align*}
    R&=\frac{\langle P\rangle}{\langle I^2\rangle}\\
    R&=\frac{\frac{p_0^2\omega^4}{12\pi\varepsilon_0c^3}}{\frac{1}{2}q_0^2\omega^2}\\
    R&=\frac{d^2\omega^2}{6\pi\varepsilon_0c^3}
\end{align*}
\mysssec{不利用假设3，计算振荡磁偶极子的电磁场。求出坡印廷矢量，证明辐射强度与利用近似3时所得到的结果相同。}
\begin{align*}
    \boldsymbol{A}&=\frac{\mu_0I_0b^2}{4}\frac{
    c\cos\left(t\omega-\frac{r\omega}{c}\right)-\sin\left(t\omega-\frac{r\omega}{c}\right)\omega r}{cr^2}\sin\theta\boldsymbol{e}_\phi\\
    \boldsymbol{E}&=\frac{\mu_0m_0\omega}{4\pi}\frac{
    \cos\left(t\omega-\frac{r\omega}{c}\right)\omega r-c\sin\left(t\omega-\frac{r\omega}{c}\right)}{cr^2}\sin\theta\boldsymbol{e}_\phi
    \end{align*}
    \begin{align*}
    \boldsymbol{B}&=\boldsymbol{\nabla}\times\boldsymbol{A}\\
    \boldsymbol{B}&=\frac{1}{r\sin\theta}\left(\frac{\partial\sin\theta  A}{\partial \theta}\right)\boldsymbol{e}_r
    +\frac{1}{r\sin\theta }\left(
    -\sin\theta\frac{\partial r A}{\partial r}
    \right)\boldsymbol{e}_\theta\\
    \boldsymbol{B}&=\frac{A}{r\tan\theta}\boldsymbol{e}_r
    +\frac{1}{r\sin\theta }\left[
    -\sin\theta\frac{\partial r \dfrac{\mu_0I_0b^2}{4}\dfrac{
    c\cos\left(t\omega-\frac{r\omega}{c}\right)-\sin\left(t\omega-\frac{r\omega}{c}\right)\omega r}{cr^2}\sin\theta}{\partial r}
    \right]\boldsymbol{e}_\theta\\
    \boldsymbol{B}&=\frac{A}{r\tan\theta}\boldsymbol{e}_r
    -\frac{1}{r}\left[
    \frac{\partial \dfrac{\mu_0I_0b^2}{4}\dfrac{
    c\cos\left(t\omega-\frac{r\omega}{c}\right)-\sin\left(t\omega-\frac{r\omega}{c}\right)\omega r}{cr}\sin\theta}{\partial r}
    \right]\boldsymbol{e}_\theta\\
    \boldsymbol{B}&=\frac{A}{r\tan\theta}\boldsymbol{e}_r
    -\frac{\sin\theta}{r}\frac{\mu_0I_0b^2}{4}\left[
    \dfrac{\sin\left(t\omega-\frac{r\omega}{c}\right)\omega r-c\cos\left(t\omega-\frac{r\omega}{c}\right)}{cr^2}
    +\dfrac{\omega^2\cos\left(t\omega-\frac{r\omega}{c}\right)}{c^2}
    \right]\boldsymbol{e}_\theta\\
    S_\phi&=\frac{1}{\mu_0}\frac{\mu_0m_0\omega}{4\pi}\frac{
    \cos\left(t\omega-\frac{r\omega}{c}\right)\omega r-c\sin\left(t\omega-\frac{r\omega}{c}\right)}{cr^2}\sin\theta\frac{A}{r\tan\theta}\\
    S_\phi&=\frac{\mu_0m_0^2\omega}{16\pi^2}\frac{
    \cos\left(t\omega-\frac{r\omega}{c}\right)\omega r-c\sin\left(t\omega-\frac{r\omega}{c}\right)}{cr^2}\frac{
    c\cos\left(t\omega-\frac{r\omega}{c}\right)-\sin\left(t\omega-\frac{r\omega}{c}\right)\omega r}{cr^2}\sin\theta\cos\theta\\
    \langle S_\phi\rangle&=0\\
    S_r&=\frac{\mu_0m_0^2\omega}{16\pi^2}\frac{
    \cos\Phi\omega r-c\sin\Phi}{cr^2}\frac{\sin^2\theta}{r}
    \dfrac{\sin\Phi\omega rc-c^2\cos\Phi
    +\omega^2r^2\cos\Phi}{c^2r^2}\\
    \langle S_r\rangle&=\langle\frac{\mu_0m_0^2\omega}{16\pi^2}\frac{\sin^2\theta}{r}
    \frac{\cos^2\Phi\omega r(-c^2
    +\omega^2r^2)
    -\sin^2\Phi\omega rc^2}
    {c^3r^4}\rangle\\
    \langle P\rangle&=\frac{\mu_0m_0^2\omega^4}{12\pi c^3}
\end{align*}
\mysssec{一个半径为$b$的绝缘的圆形环位于$xy$平面, 中心在原点. 它的电荷线密度为$\lambda=\lambda_0\sin\phi$. 现在线圈以角速度$\omega$绕$z$轴旋转. 计算辐射功率}
\begin{align*}
    \boldsymbol{d}&=\int_{0}^{2\pi}\lambda_0\sin\left(\phi+\omega t\right) \left(x\boldsymbol{e}_x+y\boldsymbol{e}_y\right)\dif\phi\\
    \boldsymbol{d}&=\int_{0}^{2\pi}\lambda_0\sin\left(\phi+\omega t\right)\left(r\cos\phi\boldsymbol{e}_x+r\sin\phi\boldsymbol{e}_y\right)\dif\phi\\
    \boldsymbol{d}&=\int_{0}^{2\pi}\lambda_0\left(\sin\phi\cos\omega t+\cos\phi\sin\omega t\right)\left(r\cos\phi\boldsymbol{e}_x+r\sin\phi\boldsymbol{e}_y\right)\dif\phi\\
    \boldsymbol{d}&=\int_{0}^{2\pi}\lambda_0r\left(\sin^2\phi\cos\omega t\boldsymbol{e}_y+\cos^2\phi\sin\omega t\boldsymbol{e}_x\right)\dif\phi\\
    \boldsymbol{d}&=\lambda_0r\pi\left(\cos\omega t\boldsymbol{e}_y+\sin\omega t\boldsymbol{e}_x\right)\\
    \ddot{\boldsymbol{d}}&=-\lambda_0r\omega^2\pi\left(\cos\omega t\boldsymbol{e}_y+\sin\omega t\boldsymbol{e}_x\right)\\
    \ddot{\boldsymbol{d}}^2&=\lambda_0^2r^2\omega^4\pi^2\\
    P&=\frac{\mu_0\ddot{\boldsymbol{d}}^2(t_0)}{6\pi c}\\
    P&=\frac{\mu_0\lambda_0^2r^2\omega^4\pi}{6c}
\end{align*}
\mysssec{作为一个电四极矩辐射模型，考虑两个相反方向振动的电偶极子，相距为$d$：\\
(a）求标势和矢势。\\
(b）求电场和磁场。\\
(c）求坡印廷矢量和辐射功率。}
(a)
\begin{align*}
    U&\cong-\frac{d_0\sin\left[\omega\left(t-\frac{r_{+}}{c}\right)\right]}{4\pi\varepsilon_0}\frac{
    \omega \cos\theta}{r_+c}
    +\frac{d_0\sin\left[\omega\left(t-\frac{r_-}{c}\right)\right]}{4\pi\varepsilon_0}\frac{\omega \cos\theta}{r_-c}\\
    U&\cong-\frac{\sin\left[\omega\left(t-\frac{r}{c}+\frac{d}{2c}\cos\theta\right)\right]}{4\pi\varepsilon_0}\frac{
    d_0\omega \cos\theta}{r_+c}
    +\frac{d_0\sin\left[\omega\left(t-\frac{r}{c}-\frac{d}{2c}\cos\theta\right)\right]}{4\pi\varepsilon_0}\frac{\omega \cos\theta}{r_-c}\\
    U&\cong\frac{d_0\omega\cos\theta}{c}\left[
    \frac{\sin\left(\omega t_0\right)
    -\cos\left(\omega t_0\right)\frac{\omega d}{2c}\cos\theta}{4\pi\varepsilon_0r_-}
    -\frac{\sin\left(\omega t_0\right)
    +\cos\left(\omega t_0\right)\frac{\omega d}{2c}\cos\theta}{4\pi\varepsilon_0r_+}\right]\\
    U&\cong\frac{d_0\omega\cos\theta}{c}\left[
    -\frac{
    \cos\left(\omega t_0\right)\omega d\cos\theta}{4\pi\varepsilon_0rc}-
    \frac{\sin\left(\omega t_0\right)
    }{4\pi\varepsilon_0r}\frac{d}{r}\cos\theta
    \right]\\
    U&\cong-\frac{d_0\omega\cos^2\theta d}{4\pi\varepsilon_0rc}\left[
    \frac{
    \cos\left(\omega t_0\right)\omega}{c}+
    \frac{\sin\left(\omega t_0\right)
    }{r}
    \right]\\
    U&\cong-\frac{d_0\omega^2\cos^2\theta d}{4\pi\varepsilon_0rc^2}
    \cos\left(\omega t_0\right)\\
    \boldsymbol{A}&\cong-\frac{\mu_0}{4\pi}\frac{d_0\omega\sin\left(t\omega-\frac{r_+\omega}{c}\right)}{r_+}\boldsymbol{e}_z
    +\frac{\mu_0}{4\pi}\frac{d_0\omega\sin\left(t\omega-\frac{r_-\omega}{c}\right)}{r_-}\boldsymbol{e}_z\\
    \boldsymbol{A}&\cong-\frac{d_0\mu_0\omega^2\cos\theta d}{4\pi rc}
    \cos\left(\omega t_0\right)\boldsymbol{e}_z
\end{align*}

(b)
\begin{align*}
    \boldsymbol{E}&=-\boldsymbol{\nabla}U-\frac{\partial \boldsymbol{A}}{\partial t}\\
    \boldsymbol{E}&=\boldsymbol{\nabla}\frac{d_0\omega^2\cos^2\theta d}{4\pi\varepsilon_0rc^2}
    \cos\left(\omega t_0\right)+\frac{\partial \frac{d_0\mu_0\omega^2\cos\theta d}{4\pi rc}
    \cos\left(\omega t_0\right)}{\partial t}\boldsymbol{e}_z\\
    \boldsymbol{E}&=\frac{d_0\omega^2d}{4\pi c}\left[\boldsymbol{\nabla}\frac{\cos^2\theta }{\varepsilon_0rc}
    \cos\left(\omega t_0\right)
    -\frac{\mu_0\omega\cos\theta
    \sin\left(\omega t_0\right)}{r}\boldsymbol{e}_z\right]\\
    \boldsymbol{\nabla}\frac{\cos^2\theta }{\varepsilon_0r}
    \cos\left(\omega t_0\right)
    &=\dfrac{\cos^2\theta\partial \dfrac{\cos\left(\omega t_0\right)}{r}
    }{\varepsilon_0\partial r}\boldsymbol{e}_r
    +\frac{\cos\left(\omega t_0\right)}{\varepsilon_0r^2}\dfrac{\partial \cos^2\theta 
    }{\partial \theta}\boldsymbol{e}_\theta\\
    \boldsymbol{\nabla}\frac{\cos^2\theta }{\varepsilon_0r}
    \cos\left(\omega t_0\right)
    &=\dfrac{\cos^2\theta\cos\left(\omega t_0\right)\partial \dfrac{1}{r}
    }{\varepsilon_0\partial r}\boldsymbol{e}_r
    +\dfrac{\cos^2\theta\dfrac{1}{r}\partial \cos\left(\omega t_0\right)
    }{\varepsilon_0\partial r}\boldsymbol{e}_r
    -\frac{2\cos\left(\omega t_0\right)\sin\theta\cos\theta}{\varepsilon_0r^2}\boldsymbol{e}_\theta\\
    \boldsymbol{\nabla}\frac{\cos^2\theta }{\varepsilon_0r}
    \cos\left(\omega t_0\right)
    &=-\dfrac{\cos^2\theta\cos\left(\omega t_0\right)
    }{\varepsilon_0r^2}\boldsymbol{e}_r
    +\dfrac{\cos^2\theta\omega\sin\left(\omega t_0\right)
    }{rc\varepsilon_0}\boldsymbol{e}_r
    -\frac{2\cos\left(\omega t_0\right)\sin\theta\cos\theta}{\varepsilon_0r^2}\boldsymbol{e}_\theta\\
    \boldsymbol{\nabla}\frac{\cos^2\theta }{\varepsilon_0r}
    \cos\left(\omega t_0\right)
    &=\dfrac{\cos^2\theta\omega\sin\left(\omega t_0\right)
    }{rc\varepsilon_0}\boldsymbol{e}_r\\
    \boldsymbol{E}&=\frac{d_0\omega^2d}{4\pi c}\left[\dfrac{\cos^2\theta\omega\sin\left(\omega t_0\right)
    }{rc^2\varepsilon_0}\boldsymbol{e}_r
    -\frac{\mu_0\omega\cos\theta
    \sin\left(\omega t_0\right)}{r}\boldsymbol{e}_z\right]\\
    \boldsymbol{E}&=\frac{d_0\omega^2d\sin\left(\omega t_0\right)\cos\theta\omega}{4\pi rc}\left[\mu_0\cos\theta
    \boldsymbol{e}_r
    -\mu_0 \left(\cos\theta{\boldsymbol{e}_r} -\sin\theta{\boldsymbol{e}_\theta} \right)\right]\\
    \boldsymbol{E}&=\frac{d_0\omega^3d\sin\left(\omega t_0\right)\cos\theta}{4\pi rc}\mu_0\sin\theta{\boldsymbol{e}_\theta}
\end{align*}
\begin{align*}
    \boldsymbol{B}&=\boldsymbol{\nabla}\times\boldsymbol{A}\\
    \boldsymbol{B}&=\boldsymbol{\nabla}\times\frac{d_0\mu_0\omega^2\cos\theta d}{4\pi rc}
    \cos\left(\omega t_0\right)\left(\cos\theta{\boldsymbol{e}_r} -\sin\theta{\boldsymbol{e}_\theta} \right)\\
    \boldsymbol{B}&=\frac{d_0\mu_0\omega^2d}{4\pi c}\boldsymbol{\nabla}\times\frac{\cos\theta }{r}
    \cos\left(\omega t_0\right)\left(-\cos\theta{\boldsymbol{e}_r} +\sin\theta{\boldsymbol{e}_\theta} \right)\\
    \boldsymbol{B}&=\frac{1}{r}\left(\frac{\partial rA_\theta}{\partial r}
    -\frac{\partial A_r}{\partial \theta}\right)\boldsymbol{e}_\phi\\
    \boldsymbol{B}&=\frac{d_0\mu_0\omega^2d}{4\pi rc}\left[\frac{\partial r\frac{\cos\theta }{r}
    \cos\left(\omega t_0\right)\sin\theta}{\partial r}
    +\frac{\partial \frac{\cos^2\theta }{r}
    \cos\left(\omega t_0\right)}{\partial \theta}\right]\boldsymbol{e}_\phi\\
    \boldsymbol{B}&=\frac{d_0\mu_0\omega^2d}{4\pi rc}\left[\frac{\cos\theta
    \sin\left(\omega t_0\right)\omega\sin\theta}{c}
    -\frac{2\sin\theta\cos\theta }{r}
    \cos\left(\omega t_0\right)\right]\boldsymbol{e}_\phi\\
    \boldsymbol{B}&=\frac{d_0\mu_0\omega^3d}{4\pi rc^2}\cos\theta
    \sin\left(\omega t_0\right)\sin\theta
    \boldsymbol{e}_\phi
\end{align*}

(c)\begin{align*}
    \boldsymbol{S}&=\frac{1}{\mu_0}\boldsymbol{E}\times\boldsymbol{B}\\
    \boldsymbol{S}&=\frac{1}{\mu_0}\frac{d_0\omega^3d\sin\left(\omega t_0\right)\cos\theta}{4\pi rc}\mu_0\sin\theta{\boldsymbol{e}_\theta}\times\frac{d_0\mu_0\omega^3d}{4\pi rc^2}\cos\theta
    \sin\left(\omega t_0\right)\sin\theta
    \boldsymbol{e}_\phi\\
    \boldsymbol{S}&=\frac{d_0^2\omega^6d^2\sin^2\left(\omega t_0\right)\cos^2\theta\mu_0}{16\pi^2 r^2c^3}\sin^2\theta
    \boldsymbol{e}_r\\
    \langle \boldsymbol{S}\rangle&=\frac{d_0^2\omega^6d^2\cos^2\theta\mu_0}{32\pi^2 r^2c^3}\sin^2\theta
    \boldsymbol{e}_r\\
    \langle P\rangle&=\int_{0}^{\pi}\frac{d_0^2\omega^6d^2\cos^2\theta\mu_0}{32\pi^2 r^2c^3}\sin^3\theta 4\pi r^2\dif \theta\\
    \langle P\rangle&=-\int_{0}^{\pi}\frac{d_0^2\omega^6d^2\cos^2\theta\mu_0}{8\pi c^3}\left(1-\cos^2\theta\right)\dif \cos\theta\\
    \langle P\rangle&=\frac{d_0^2\omega^6d^2\mu_0}{30\pi c^3}
\end{align*}
\subsection{点电荷}
\subsubsection{点电荷的辐射功率}

考虑电荷 $q$ 在任意时刻的位置为 $\boldsymbol{w}(t)$，
速度和加速度分别为
\begin{align*}
    \boldsymbol{v}=\dot{\boldsymbol{w}},
    \qquad
    \boldsymbol{a}=\dot{\boldsymbol{v}}.
\end{align*}

我们已经得到任意运动点电荷的Liénard--Wiechert 场。式 \ref{eq:LW_E_radiation} 中的第一项称为
\textbf{速度场}，第二项称为
\textbf{加速场}。

\begin{align*}
    \boldsymbol S
    &=
    \frac{1}{\mu_0}
    \boldsymbol E\times\boldsymbol B\\
    \boldsymbol B
    &=
    \frac{1}{c}{\boldsymbol{e}_r}\times\boldsymbol E\\
    \boldsymbol S
    &=
    \frac{1}{\mu_0c}
    \boldsymbol E\times
    \left(
        \boldsymbol{e}_r\times\boldsymbol E
    \right)
    \\
    &=
    \frac{1}{\mu_0c}
    \left[
        E^2\boldsymbol{e}_r
        -
        (\boldsymbol{e}_r\cdot\boldsymbol E)\boldsymbol E
    \right]
\end{align*}


在远区 $r\rightarrow\infty$ 时，
速度场中的电场满足
\begin{equation}
    E_{\mathrm{vel}}\sim\frac{1}{r^2},
\end{equation}
而加速场满足
\begin{equation}
    E_{\mathrm{acc}}\sim\frac{1}{r}.
\end{equation}

由于坡印廷矢量满足
\begin{equation}
    S\sim E^2,
\end{equation}
因此只有加速场产生
\begin{equation}
    S_{\mathrm{rad}}\sim\frac{1}{r^2}
\end{equation}
的能流。

这是辐射场的重要特征：
虽然辐射场本身随 $1/r$ 衰减，
但通过半径为 $r$ 的球面的总功率在 $r\rightarrow\infty$ 时保持有限。因此只需保留式 \ref{eq:LW_E_radiation}
中的加速场：
\begin{equation}
    \boxed{
    \boldsymbol E_{\mathrm{rad}}
    =
    \frac{q}{4\pi\epsilon_0}
    \frac{
        \boldsymbol r\times
        (\boldsymbol u\times\boldsymbol a)
    }
    {(\boldsymbol r\cdot\boldsymbol u)^3}
    }
    \label{eq:Erad_general}
\end{equation}

在辐射区中，
\begin{equation}
    \boldsymbol E_{\mathrm{rad}}
    \perp\boldsymbol{e}_r,
\end{equation}
因此
\begin{equation}
    \boldsymbol S_{\mathrm{rad}}
    =
    \frac{1}{\mu_0c}
    E_{\mathrm{rad}}^2\boldsymbol{e}_r.
    \label{eq:Srad_general}
\end{equation}

首先考虑
\begin{equation}
    v\ll c.
\end{equation}

此时可以近似取
\begin{equation}
    \boldsymbol u\simeq c\boldsymbol{e}_r.
\end{equation}

代入式 \eqref{eq:Erad_general}，
得到
\begin{align}
    \boldsymbol E_{\mathrm{rad}}
    &\simeq
    \frac{q}{4\pi\epsilon_0}
    \frac{
        c\boldsymbol{e}_r
        \times
        \left[
            c\boldsymbol{e}_r\times\boldsymbol a
        \right]
    }
    {(cr)^3}
    \\
    &=
    \frac{q}{4\pi\epsilon_0c^2r}
    \left[
        \boldsymbol{e}_r
        \times
        \left(
            \boldsymbol{e}_r\times\boldsymbol a
        \right)
    \right].
\end{align}

利用向量三重积
\begin{equation}
    \boldsymbol A\times
    (\boldsymbol B\times\boldsymbol C)
    =
    \boldsymbol B(\boldsymbol A\cdot\boldsymbol C)
    -
    \boldsymbol C(\boldsymbol A\cdot\boldsymbol B),
\end{equation}
可以写成
\begin{equation}
    \boxed{
    \boldsymbol E_{\mathrm{rad}}
    =
    \frac{q}{4\pi\epsilon_0c^2r}
    \left[
        (\boldsymbol{e}_r\cdot\boldsymbol a)\boldsymbol{e}_r
        -
        \boldsymbol a
    \right]
    }.
\end{equation}

设 $\theta$ 为观察方向 $\boldsymbol{e}_r$
与加速度 $\boldsymbol a$ 之间的夹角，则
\begin{equation}
    E_{\mathrm{rad}}
    =
    \frac{qa\sin\theta}
    {4\pi\epsilon_0c^2r}.
\end{equation}

因此
\begin{equation}
    S_{\mathrm{rad}}
    =
    \frac{1}{\mu_0c}
    E_{\mathrm{rad}}^2
    =
    \frac{\mu_0q^2a^2}
    {16\pi^2cr^2}
    \sin^2\theta.
\end{equation}

于是单位立体角内的辐射功率为
\begin{equation}
    \boxed{
    \frac{dP}{d\Omega}
    =
    r^2S_{\mathrm{rad}}
    =
    \frac{\mu_0q^2a^2}
    {16\pi^2c}
    \sin^2\theta
    }.
    \label{eq:Larmor_angular}
\end{equation}

这说明点电荷的辐射具有偶极辐射的角分布。

在加速度方向
\begin{equation}
    \theta=0,\pi
\end{equation}
处没有辐射，而在垂直于加速度的方向
\begin{equation}
    \theta=\frac{\pi}{2}
\end{equation}
处辐射最强。

对所有方向积分：
\begin{equation}
    P
    =
    \int
    \frac{dP}{d\Omega}\,d\Omega.
\end{equation}

代入式 \eqref{eq:Larmor_angular}，
\begin{align}
    P
    &=
    \frac{\mu_0q^2a^2}{16\pi^2c}
    \int_0^{2\pi}d\phi
    \int_0^\pi
    \sin^3\theta\,d\theta
    \\
    &=
    \frac{\mu_0q^2a^2}{16\pi^2c}
    (2\pi)
    \left(\frac{4}{3}\right).
\end{align}

因此得到
\begin{equation}
    \boxed{
    P
    =
    \frac{\mu_0q^2a^2}{6\pi c}
    }.
    \label{eq:Larmor}
\end{equation}

利用
\begin{equation}
    \mu_0\epsilon_0c^2=1,
\end{equation}
也可以写成
\begin{equation}
    \boxed{
    P
    =
    \frac{q^2a^2}
    {6\pi\epsilon_0c^3}
    }.
\end{equation}

这就是
\textbf{Larmor 公式}。

由几何关系可以得到
\begin{equation}
    \frac{dt}{dt_r}
    =
    \frac{\boldsymbol{e}_r\cdot\boldsymbol u}{c},
\end{equation}
因此
\begin{equation}
    \frac{dW}{dt_r}
    =
    \frac{\boldsymbol{e}_r\cdot\boldsymbol u}{c}
    \frac{dW}{dt}.
\end{equation}

另一方面，辐射场给出的单位立体角辐射功率为
\begin{equation}
    \boxed{
    \frac{dP}{d\Omega}
    =
    \frac{q^2}{16\pi^2\epsilon_0}
    \frac{
        \left|
        \boldsymbol{e}_r
        \times
        (\boldsymbol u\times\boldsymbol a)
        \right|^2
    }
    {(\boldsymbol{e}_r\cdot\boldsymbol u)^5}
    }.
    \label{eq:Lienard_angular}
\end{equation}

这就是相对论情况下的辐射角分布。

与非相对论结果相比，
这里最重要的变化是分母中出现了
\begin{equation}
    (\boldsymbol{e}_r\cdot\boldsymbol u)^5,
\end{equation}
使得辐射在高速运动时产生强烈的方向性。

对式 \eqref{eq:Lienard_angular}
在整个立体角上积分，可以得到一般速度下的总辐射功率：
\begin{equation}
    \boxed{
    P
    =
    \frac{\mu_0q^2\gamma^6}{6\pi c}
    \left(
        a^2
        -
        \frac{
            |\boldsymbol v\times\boldsymbol a|^2
        }{c^2}
    \right)
    }
    \label{eq:Lienard}
\end{equation}
其中
\begin{equation}
    \gamma
    =
    \frac{1}
    {\sqrt{1-v^2/c^2}}.
\end{equation}

这就是
\textbf{Liénard 公式}。

利用
\begin{equation}
    |\boldsymbol v\times\boldsymbol a|^2
    =
    v^2a^2-(\boldsymbol v\cdot\boldsymbol a)^2,
\end{equation}
还可以将其写成
\begin{equation}
    \boxed{
    P
    =
    \frac{\mu_0q^2\gamma^4}{6\pi c}
    \left[
        a^2
        +
        \frac{\gamma^2}{c^2}
        (\boldsymbol v\cdot\boldsymbol a)^2
    \right]
    }.
\end{equation}

若
\begin{equation}
    \boldsymbol v\perp\boldsymbol a,
\end{equation}
则
\begin{align}
    P
    &=
    \frac{\mu_0q^2\gamma^6}{6\pi c}
    a^2
    \left(
        1-\frac{v^2}{c^2}
    \right)
    \\
    &=
    \boxed{
    \frac{\mu_0q^2\gamma^4a^2}
    {6\pi c}
    }.
\end{align}

例如带电粒子在磁场中做圆周运动时，
加速度始终垂直于速度，因此其辐射功率具有上述形式。

若
\begin{equation}
    \boldsymbol v\parallel\boldsymbol a,
\end{equation}
则
\begin{equation}
    \boxed{
    P=\frac{\mu_0q^2\gamma^6a^2}{6\pi c}
    }.
\end{equation}

\subsubsection{辐射反作用}
假设在周期运动下
\begin{align}
    \int_{0}^{T}\boldsymbol{F}\cdot\boldsymbol{v}\dif t&=-\frac{\mu_0q^2}{6\pi c}\int_{0}^{T}a^2\dif t\\
    \int_{0}^{T}\boldsymbol{F}\cdot\boldsymbol{v}\dif t&=-\frac{\mu_0q^2}{6\pi c}\int_{0}^{T}\frac{\dif\boldsymbol{v}}{\dif t}\cdot\frac{\dif\boldsymbol{v}}{\dif t}\dif t\nonumber\\
    \int_{0}^{T}\boldsymbol{F}\cdot\boldsymbol{v}\dif t&=-\frac{\mu_0q^2}{6\pi c}\left.\boldsymbol{v}\cdot\frac{\dif\boldsymbol{v}}{\dif t}\right|_{0}^{T}
    +\frac{\mu_0q^2}{6\pi c}\int_{0}^{T}\boldsymbol{v}\cdot\frac{\dif^2\boldsymbol{v}}{\dif t^2}\dif t\nonumber\\
    \int_{0}^{T}\boldsymbol{F}\cdot\boldsymbol{v}\dif t&=\int_{0}^{T}\boldsymbol{v}\cdot\dif t\nonumber\\
    \int_{0}^{T}\left(\boldsymbol{F}-\frac{\mu_0q^2}{6\pi c}\frac{\dif^2\boldsymbol{v}}{\dif t^2}\right)\cdot\boldsymbol{v}\dif t&=0
\end{align}
上式仅为猜测。关于$\boldsymbol{F}$垂直于的分量得不出任何结论，仅告诉你平行分量在一个特别的时间间隔内的平均。
\subsubsection{辐射反作用的物理基础}
点电荷的场在粒子所在的地方为无限大，很难明白如何计算它们施加的力。为避免这个问题，让我们考虑扩展电荷分布的情形，这种情形下场在任何地方都是有限值，最后我们考虑电荷的尺寸趋于零时的极限。一般地，一部分（A）对另一部分（B）的电磁力与B作用在A上的力大小不同，方向不相反。如果把分布分为无限小，则每一对都存在着这种不平衡力，对所有的不平衡力求和，结果是电荷作用在自身上的一个净的作用力。正是这种自作用力导致在粒子内部牛顿第三定律被打破，这是辐射反作用的作用。

模型：一个带总电荷为$2q$的“哑铃”，电荷平均分布在两个相距为$d$的球上。

假设哑铃沿$x$方向运动，在推迟时刻（瞬时）静止。（A）在（B）处产生的电场是
$$E_{AB}=\frac{q}{4\pi\varepsilon_0}$$
\mysssec{假设一个初速度为$v_0$的电子以恒定加速度减速至速度为零。它辐射损失的能量占它初始动能的比值是多少？(其余的能量被为保持它以一个恒定减速运动的其他机制吸收掉了。）假设$v_0<<c$，所以拉莫尔公式适用。}
\begin{align*}
    t&=\frac{v_0}{a}\\
    W&=Pt\\
    W&=\frac{q^2a^2}
    {6\pi\epsilon_0c^3}\frac{v_0}{a}\\
    \frac{2W}{mv_0^2}&=\frac{q^2a}
    {3\pi\epsilon_0c^3mv_0}
\end{align*}
\mysssec{(a)带电荷为$q$的粒子，以速度$v$沿半径为$R$的圆周作匀速圆周运动。为了保持这个运动，当然需要
提供一个向心力。为了抵消辐射反作用，必须施加的附加力$F$是多少？这个附加力传递的功率$P_e$是多少？\\
(b)对作简谐振动的粒子重复(a)的计算。\\
(c)考虑一个粒子自由下落的情况（加速度为常数$g$）。辐射反作用力是多少？}
(a)\begin{align*}
    \boldsymbol F&=\frac{\mu_0q^2}{6\pi c}\frac{\dif^2\boldsymbol{v}}{\dif t^2}\\
    \boldsymbol F&=-\frac{\mu_0q^2}{6\pi c}\frac{v^2}{r^2}\boldsymbol{v}\\
    P_e&=\frac{\mu_0q^2}{6\pi c}\frac{v^4}{r^2}
\end{align*}

(b)\begin{align*}
    \boldsymbol F&=\frac{\mu_0q^2}{6\pi c}\frac{\dif^2\boldsymbol{v}}{\dif t^2}\\
    \boldsymbol F&=-\frac{\mu_0q^2}{6\pi c}\frac{v^2}{r^2}\boldsymbol{v}\\
    P_e&=\frac{\mu_0q^2}{6\pi c}\frac{v^4}{r^2}
\end{align*}
\mysssec{考虑辐射反作用力，带电粒子的牛顿第二定律变为
$$a=\tau\dot{a}+\frac{F}{m}$$
（a）与不带电粒子不同（$a=F/m$），加速度必须是时间的连续函数，即使力突然变化。（从物理上看，辐射反作用阻碍a的快速变化。）通过对上面运动方程从（$t-\varepsilon$）到（$t+\varepsilon$）积分并取极
限$\varepsilon\to 0$，证明a在任何时刻都是连续的。\\
（b）一个粒子在$t=0$受到一恒定的力$F$，持续时间为$T$。求在下面三个时间内最一般的运动方程的解：
(i)$t<0$;(ii) $0<t<T$;(iii)$t>T$。}
(a)
\begin{align*}
    \lim_{\varepsilon\to 0}\int_{t-\varepsilon}^{t+\varepsilon}a\dif t&=\lim_{\varepsilon\to 0}\int_{t-\varepsilon}^{t+\varepsilon}\tau\dot{a}+\frac{F}{m}\dif t\\
    a_{t+\varepsilon}-a_{t-\varepsilon}&=0
\end{align*}

(b)$t<0$\begin{align*}
    a&=\tau\frac{\dif a}{\dif t}\\
    \dif t&=\tau\frac{\dif a}{a}\\
    t&=\tau\ln a+C_0\\
    a&=a_0e^{t/\tau}
\end{align*}

$0<t<T$\begin{align*}
    a-\frac{F}{m}&=\tau\frac{\dif a}{\dif t}\\
    \dif t&=\tau\frac{\dif a}{a-\frac{F}{m}}\\
    t&=\tau\ln \left(a-\frac{F}{m}\right)+C\\
    a&=\left(a_0-\frac{F}{m}\right)e^{t/\tau}+\frac{F}{m}
\end{align*}

$t>T:a=\left[\left(a_0-\frac{F}{m}\right)e^{T/\tau}+\frac{F}{m}\right]e^{t/\tau}$